\documentclass[12pt,a4paper]{article}
\usepackage[utf8]{inputenc}
\usepackage[french]{babel}
\usepackage[T1]{fontenc}
\usepackage{geometry}
\usepackage{enumitem}
\usepackage{titlesec}
\usepackage{hyperref}
\usepackage{xcolor}

\geometry{top=2.5cm, bottom=2.5cm, left=2.5cm, right=2.5cm}

\titleformat{\section}
{\normalfont\Large\bfseries\color{blue!70!black}}{\thesection}{1em}{}

\titleformat{\subsection}
{\normalfont\large\bfseries\color{blue!50!black}}{\thesubsection}{1em}{}

\title{\textbf{Cahier des Charges\\Système de Gestion de Salles}}
\author{}
\date{Version 2.0 - \today}

\begin{document}

\maketitle
\tableofcontents
\newpage

\section{Introduction}

Le présent cahier des charges définit les spécifications pour le développement d'un système de gestion des salles de l'école. Ce système a pour objectif d'optimiser la gestion des salles en fonction des besoins des filières, des professeurs et des membres de clubs, et de permettre une gestion efficace des réservations et de la disponibilité des espaces.

\textbf{⚠️ Version 2.0 - Architecture Unifiée :} Cette version adopte un système de gestion des utilisateurs unifié avec une table unique et un système de rôles, simplifiant l'architecture et centralisant la gestion des inscriptions.

\section{Objectifs du Système}

Le système doit permettre :

\begin{itemize}
    \item \textbf{Aux membres d'administration} de gérer tous les utilisateurs (approuver les inscriptions, gérer les filières, les salles, traiter les réclamations).
    \item \textbf{Aux coordinateurs de filière} d'ajouter des filières, de définir leurs effectifs selon les cycles (Préparatoire 2 ans, Ingénieur 3 ans), les charges horaires des matières (cours, TP, TD avec groupes) et de créer les emplois du temps manuellement.
    \item \textbf{Aux professeurs} de consulter leur emploi du temps, libérer des salles (exceptionnellement ou définitivement), demander des réservations pour des séances supplémentaires, et faire des réclamations pour des équipements.
    \item \textbf{Aux membres de clubs} de s'inscrire, réserver des salles pour leurs activités (sauf salles TP), et libérer des salles après utilisation.
    \item D'optimiser la gestion des réservations selon une politique de priorité, avec priorité aux professeurs (100) et FIFO au sein de chaque catégorie d'acteurs.
\end{itemize}

\section{Acteurs et Rôles}

\subsection{Membre d'Administration}
\textbf{Rôle central} du système, créé automatiquement au démarrage de l'application.

\textbf{Responsabilités :}
\begin{itemize}
    \item \textbf{Gestion des inscriptions :} Recevoir les notifications (Email + WhatsApp) pour chaque nouvelle inscription, approuver ou refuser les inscriptions des utilisateurs.
    \item \textbf{Attribution des rôles :} Lors de l'approbation, attribuer la filière pour les Professeurs et Coordinateurs, ou saisir le nom du club pour les Membres de Club.
    \item \textbf{Gestion des utilisateurs :} Consulter, activer ou désactiver les comptes utilisateurs.
    \item \textbf{Gestion des salles :} CRUD complet des salles (nom, localisation, capacité, type).
    \item \textbf{Gestion des filières :} Définir les effectifs des étudiants par filière et année.
    \item \textbf{Traitement des réclamations :} Consulter et traiter les réclamations d'équipements.
    \item \textbf{Consultation des réservations :} Vue globale de toutes les réservations.
\end{itemize}

\subsection{Coordinateur de Filière}
Responsable de la gestion de sa filière et de la création manuelle des emplois du temps. \textbf{Remarque :} Il y a un coordinateur par filière (non global). Un coordinateur est également un professeur (cumul de rôles possible).

\textbf{Responsabilités spécifiques :}
\begin{itemize}
    \item Définir les charges horaires des matières (Cours, TP, TD) par filière
    \item Créer manuellement les emplois du temps via interface calendrier interactive avec onglets par année (ex: DLA1, DLA2, DLA3)
    \item Sélectionner le groupe (Groupe 1 ou Groupe 2) pour les séances de TP et TD
    \item Toutes les fonctionnalités du professeur
\end{itemize}

\subsection{Professeur}
Enseignant de l'école, s'inscrit via le formulaire commun puis est approuvé par le Membre d'Administration.

\textbf{Responsabilités :}
\begin{itemize}
    \item Consulter son emploi du temps personnel
    \item Libérer des salles (exceptionnellement pour une séance ou définitivement)
    \item Demander des réservations pour des séances supplémentaires (hors emploi du temps)
    \item Faire des réclamations pour des équipements
    \item Consulter la liste des salles disponibles
\end{itemize}

\subsection{Membre de Club}
Étudiant membre d'un club, s'inscrit via le formulaire commun en indiquant le nom de son club (texte libre), puis est approuvé par le Membre d'Administration.

\textbf{Responsabilités :}
\begin{itemize}
    \item Réserver des salles pour les activités du club (toutes sauf les salles TP)
    \item Libérer des salles après utilisation
    \item Faire des réclamations pour des équipements
    \item Consulter la liste des salles disponibles
\end{itemize}

\textbf{Remarque :} Les réservations des membres de clubs ont une priorité inférieure (50) à celles des professeurs (100).

\section{Fonctionnalités Détaillées}

\subsection{Fonctionnalités du Membre d'Administration}

\subsubsection{Gestion des Inscriptions}
\begin{itemize}
    \item \textbf{Workflow d'inscription :}
    \begin{enumerate}
        \item Utilisateur s'inscrit via formulaire commun (nom, prénom, email, mot de passe, rôle souhaité)
        \item Notification immédiate (Email + WhatsApp) envoyée au Membre d'Administration
        \item Inscription apparaît dans liste "En attente"
        \item Membre d'Administration accepte ou refuse
        \item Si acceptation :
        \begin{itemize}
            \item Pour Professeur/Coordinateur : Sélectionner la filière
            \item Pour Membre de Club : Saisir le nom du club (texte libre)
        \end{itemize}
        \item Statut devient "Actif" et utilisateur peut se connecter
    \end{enumerate}
    \item Consultation de toutes les inscriptions avec filtrage par statut
    \item Gestion des refus avec notification automatique
\end{itemize}

\subsubsection{Gestion des Utilisateurs}
\begin{itemize}
    \item Liste complète de tous les utilisateurs avec leurs rôles
    \item Filtrage par rôle (Admin, Coordinateur, Professeur, Membre Club)
    \item Activation/Désactivation des comptes
    \item Modification des informations utilisateur
    \item Historique des actions par utilisateur
\end{itemize}

\subsubsection{Gestion des Salles}
\begin{itemize}
    \item CRUD complet des salles
    \item Définir le type (COURS, TP, TD) qui détermine l'accessibilité pour les clubs
    \item Gérer la capacité, les équipements, la localisation
\end{itemize}

\subsubsection{Gestion des Filières}
\begin{itemize}
    \item Définir les cycles : Préparatoire (CP1, CP2) et Ingénieur (GI1-GI2-GI3, DLA1-DLA2-DLA3, etc.)
    \item Spécifier l'effectif par année de filière
    \item Attribuer un coordinateur à chaque filière
\end{itemize}

\subsubsection{Traitement des Réclamations}
\begin{itemize}
    \item Recevoir notification pour chaque nouvelle réclamation
    \item Consulter toutes les réclamations avec filtrage (urgence, statut, salle)
    \item Traiter les réclamations (passer de "En attente" à "Traitée")
    \item Ajouter des commentaires sur le traitement
    \item Notification automatique au demandeur lors du traitement
\end{itemize}

\subsubsection{Consultation des Réservations}
\begin{itemize}
    \item Vue globale de toutes les réservations (professeurs et clubs)
    \item Filtrage par date, salle, statut, demandeur
    \item Statistiques d'utilisation des salles
\end{itemize}

\subsection{Fonctionnalités du Coordinateur de Filière}

\subsubsection{Gestion de la Charge Horaire}
\begin{itemize}
    \item Définir les charges horaires des matières par filière pour les différentes catégories de séances (Cours, TP, TD)
    \item Chaque matière peut avoir des heures différentes selon le type : heures de cours, heures de TP, heures de TD
\end{itemize}

\subsubsection{Création des Emplois du Temps}
\begin{itemize}
    \item \textbf{Interface avec onglets par année :} Après sélection de la filière (ex: DLA), le coordinateur voit 3 onglets : DLA1, DLA2, DLA3
    
    \item \textbf{Création manuelle via calendrier interactif :} Grille Jours (colonnes) × Heures (lignes) avec bouton [+] sur chaque créneau
    
    \item \textbf{Modal de sélection :} Au clic sur [+], un modal s'affiche pour saisir :
    \begin{itemize}
        \item Matière (dropdown)
        \item Type de séance (Cours / TP / TD)
        \item \textbf{Si TP ou TD :} Sélection obligatoire du groupe (Groupe 1 ou Groupe 2)
        \item Salle (filtrée automatiquement par type et capacité >= effectif)
        \item Professeur (dropdown)
    \end{itemize}
    
    \item \textbf{Périodes de 7 semaines :} Les emplois du temps sont créés pour des périodes spécifiques de 7 semaines avec dates précises (exemple : 1 sept - 19 oct)
    
    \item \textbf{Génération automatique des séances :} Le système génère 7 séances individuelles (une par semaine) avec dates calculées
    
    \item Le système peut changer l'emploi du temps à chaque nouvelle période de 7 semaines
\end{itemize}

\subsection{Fonctionnalités du Professeur}

\subsubsection{Consultation de l'Emploi du Temps}
\begin{itemize}
    \item Consulter son emploi du temps personnel avec toutes ses séances
    \item Filtrer par date, par matière
    \item Vue calendrier avec indication des séances planifiées
\end{itemize}

\subsubsection{Libération de Salle}
\begin{itemize}
    \item \textbf{Libération définitive :} Lorsqu'un cours a terminé sa charge horaire, le professeur peut libérer la salle de façon permanente pour cette matière, libérant ainsi toutes les séances futures restantes
    \item \textbf{Libération exceptionnelle :} Le professeur peut libérer UNE séance spécifique pour une date future précise uniquement (exemple : libérer la séance du lundi 15 octobre seulement, les autres lundis restent bloqués). Les séances passées ne peuvent pas être libérées
    \item Les notifications de libération sont envoyées au coordinateur de filière et au membre d'administration par Email + WhatsApp
\end{itemize}

\subsubsection{Demande de Réservation de Salle}
\begin{itemize}
    \item \textbf{Réservations pour séances supplémentaires :} Les professeurs utilisent ce module uniquement pour des séances qui ne sont pas dans l'emploi du temps fixe (exemples : cours de rattrapage, séances de révision).
    \item Faire une demande de réservation en précisant :
    \begin{itemize}
        \item Nature de la séance (cours, TP, TD)
        \item Date et créneau horaire souhaité
        \item Filière associée
    \end{itemize}
    \item \textbf{Validation automatique :} La réservation est automatiquement acceptée si la salle est disponible (vérification dans l'emploi du temps et les autres réservations).
    \item Si la salle n'est pas disponible, le système propose automatiquement des alternatives (autres créneaux ou autres salles).
    \item Confirmation envoyée par email et WhatsApp.
\end{itemize}

\subsubsection{Réclamation d'Équipements}
\begin{itemize}
    \item Faire une réclamation pour des équipements (exemple : datashow, micro, tableau) pour une séance ou une salle
    \item Préciser le type d'équipement, la description du problème, et le niveau d'urgence (Faible, Moyen, Urgent)
    \item Suivre l'état de la réclamation : "En attente" ou "Traitée"
    \item Notification immédiate (Email + WhatsApp) envoyée au membre d'administration
    \item Recevoir une notification par Email + WhatsApp lors du traitement de la réclamation
\end{itemize}

\subsection{Fonctionnalités du Membre de Club}

\subsubsection{Inscription}
\begin{itemize}
    \item S'inscrire sur la plateforme web via le formulaire commun
    \item Indiquer nom, prénom, email, mot de passe et sélectionner rôle "Membre de Club"
    \item Statut initial : "En attente"
    \item Notification (Email + WhatsApp) envoyée au membre d'administration
    \item Membre d'administration approuve et saisit le nom du club (texte libre : ex. "Club Robotique", "Club Sport")
    \item Après approbation, le compte devient "Actif" et le membre peut se connecter
\end{itemize}

\subsubsection{Réservation de Salle}
\begin{itemize}
    \item Consulter les salles disponibles : toutes les salles sauf les salles de type TP.
    \item Faire une demande de réservation en précisant :
    \begin{itemize}
        \item Date et créneau horaire souhaité
        \item Nom du club
        \item Motif de la réservation (activité)
    \end{itemize}
    \item \textbf{Système de priorité :} Les réservations des professeurs (priorité 100) sont traitées avant celles des clubs (priorité 50). Au sein des clubs, la politique FIFO (premier arrivé, premier servi) s'applique.
    \item \textbf{Validation automatique :} Si la salle est disponible (pas dans l'emploi du temps et pas déjà réservée), la réservation est automatiquement acceptée.
    \item Recevoir une confirmation par email et WhatsApp en cas de réservation acceptée.
    \item Recevoir une proposition alternative par email et WhatsApp si la salle n'est pas disponible.
\end{itemize}

\subsubsection{Libération de Salle}
\begin{itemize}
    \item Libérer une salle réservée en cas d'annulation ou de fin d'activité
    \item Notification automatique (Email + WhatsApp) envoyée au membre d'administration
\end{itemize}

\subsubsection{Réclamation d'Équipements}
\begin{itemize}
    \item Même fonctionnalité que les professeurs
    \item Faire une réclamation pour des équipements défectueux ou manquants
    \item Préciser type, description, urgence
    \item Notifications Email + WhatsApp lors de création et traitement
\end{itemize}

\subsection{Fonctionnalités Générales}

\subsubsection{Gestion des Conflits de Réservation}
\begin{itemize}
    \item Si une salle n'est pas disponible pour une demande de réservation, le système propose automatiquement d'autres créneaux ou salles disponibles.
    \item \textbf{Important :} Si une salle est bloquée par l'emploi du temps fixe (séances planifiées), elle ne peut pas être réservée par les professeurs ni les clubs, sauf si la séance a été libérée exceptionnellement.
\end{itemize}

\subsubsection{Politique de Réservation et Priorités}
\begin{itemize}
    \item \textbf{Priorité des professeurs :} Les professeurs ont toujours la priorité (100) sur les membres de clubs (50) pour les réservations de salles supplémentaires.
    \item Au sein de chaque catégorie d'acteurs (professeurs entre eux, membres de clubs entre eux), les demandes de réservation sont traitées selon une politique FIFO (premier arrivé, premier servi).
    \item \textbf{Vérification de disponibilité double :} Le système vérifie à la fois l'emploi du temps fixe (table seances) et les réservations dynamiques (table reservations) avant d'accepter une demande.
    \item Les libérations exceptionnelles créent des "trous" temporaires dans l'emploi du temps fixe, permettant les réservations pour ces créneaux spécifiques.
\end{itemize}

\section{Spécifications Techniques}

\subsection{Architecture de l'Application}

\subsubsection{Application Web}
\begin{itemize}
    \item Développée en Java EE classique (Servlets, JSP).
    \item Utilisation du pattern MVC pour la structuration du code.
    \item Interface utilisateur développée avec JSP, HTML, et Tailwind CSS pour le design moderne et responsive.
    \item Accessible par tous les utilisateurs (coordinateurs, responsables de salles, professeurs, membres de clubs).
    \item Permet l'interaction avec les différentes fonctionnalités selon le profil de l'utilisateur.
\end{itemize}

\subsubsection{Conteneurisation avec Docker}
\begin{itemize}
    \item L'application est conteneurisée avec Docker pour faciliter le déploiement et la portabilité.
    \item Utilisation de Docker Compose pour orchestrer les conteneurs (application Tomcat, base de données PostgreSQL).
\end{itemize}

\subsection{Base de Données}
\begin{itemize}
    \item \textbf{Système de Gestion :} PostgreSQL
    \item \textbf{Gestion des Informations :} Stocker les informations sur les utilisateurs, les filières, les charges horaires, les emplois du temps, les salles, les réservations, les historiques de libération, et les réclamations d'équipements.
    \item \textbf{⚠️ Architecture Unifiée :}
    \begin{itemize}
        \item \texttt{users} : \textbf{TABLE UNIQUE} pour tous les utilisateurs avec système de rôles
        \begin{itemize}
            \item Champs : id, nom, prenom, email, password (hashé), role (ADMIN, COORDINATEUR, PROFESSEUR, MEMBRE\_CLUB)
            \item statut (EN\_ATTENTE, ACTIF, INACTIF, REFUSE)
            \item filiere\_id (NULL pour ADMIN et MEMBRE\_CLUB)
            \item nom\_club (texte libre, NULL sauf pour MEMBRE\_CLUB)
            \item date\_inscription, date\_approbation, approved\_by
        \end{itemize}
    \end{itemize}
    \item \textbf{Structure clé :}
    \begin{itemize}
        \item \texttt{filieres} : Code, nom, cycle (PREPARATOIRE/INGENIEUR), annee (1-2-3), effectif, coordinateur\_id (FK vers users)
        \item \texttt{matieres} : Code, nom, filiere\_id, heures\_cours, heures\_td, heures\_tp, professeur\_id (FK vers users), semestre
        \item \texttt{salles} : Numero, nom, type (COURS/TP/TD), capacite, equipements, batiment, etage, disponible
        \item \texttt{emplois\_du\_temps} : Modèles avec périodes 7 semaines, filiere\_id, matiere\_id, salle\_id, professeur\_id, jour, heure\_debut, heure\_fin, type\_seance, groupe (1 ou 2 pour TP/TD), date\_debut\_periode, date\_fin\_periode
        \item \texttt{seances} : Séances individuelles générées (date spécifique, emploi\_du\_temps\_id, salle\_id, statut: PLANIFIEE/LIBEREE\_EXCEPTIONNELLE/LIBEREE\_DEFINITIVE/ANNULEE/TERMINEE)
        \item \texttt{liberations\_exceptionnelles} : seance\_id, professeur\_id (FK vers users), motif, date\_liberation, notified\_coordinateur, notified\_admin
        \item \texttt{reservations} : salle\_id, date, heure\_debut, heure\_fin, user\_id (FK vers users), role (PROFESSEUR/MEMBRE\_CLUB), matiere\_id, motif, priorite (100 prof / 50 club), statut (EN\_ATTENTE/ACCEPTEE/REFUSEE), date\_creation
        \item \texttt{reclamations} : user\_id (FK vers users), salle\_id, seance\_id (optionnel), type\_equipement, description, urgence (FAIBLE/MOYEN/URGENT), statut (EN\_ATTENTE/TRAITEE), date\_creation, date\_traitement, commentaire\_admin
    \end{itemize}
    \item \textbf{Politique de Gestion des Conflits :} Intégrer une gestion FIFO (date\_creation) dans la base de données pour organiser les demandes de réservation au sein de chaque catégorie d'acteurs (priorité) et garantir une équité de traitement.
\end{itemize}

\subsection{Notifications}
\begin{itemize}
    \item \textbf{Canaux de notification :} Email et WhatsApp uniquement (notifications doubles pour chaque événement)
    \item \textbf{Types de notifications :}
    \begin{itemize}
        \item \textbf{Nouvelle inscription} → Membre d'Administration (Email + WhatsApp)
        \item \textbf{Approbation/Refus inscription} → Utilisateur concerné (Email + WhatsApp)
        \item \textbf{Nouvelle réclamation} → Membre d'Administration (Email + WhatsApp)
        \item \textbf{Traitement réclamation} → Demandeur (Email + WhatsApp)
        \item \textbf{Confirmation réservation} → Demandeur (Email + WhatsApp)
        \item \textbf{Refus réservation + alternatives} → Demandeur (Email + WhatsApp)
        \item \textbf{Libération salle} → Coordinateur + Membre d'Administration (Email + WhatsApp)
    \end{itemize}
    \item \textbf{Destinataires :} Membre d'Administration, Coordinateurs, Professeurs, Membres de clubs selon les cas
    \item \textbf{Intégration :} JavaMail API pour emails, Twilio API ou WhatsApp Business API pour messages WhatsApp
    \item \textbf{Service unifié :} NotificationService avec méthodes sendEmail() et sendWhatsApp()
\end{itemize}

\subsection{Sécurité et Permissions}
\begin{itemize}
    \item \textbf{Accès sécurisé :} Chaque utilisateur se connecte avec un compte sécurisé (email + mot de passe hashé)
    \item \textbf{Système de rôles :} Les permissions sont assignées en fonction du rôle (ADMIN, COORDINATEUR, PROFESSEUR, MEMBRE\_CLUB)
    \item \textbf{Workflow d'approbation :} Toute inscription nécessite l'approbation du Membre d'Administration avant activation
    \item \textbf{Compte Admin :} Créé automatiquement au démarrage de l'application (seed data)
    \item \textbf{Journalisation :} Toutes les modifications (inscription, approbation, ajout/suppression de filière, salle, libération, réservation, réclamation) sont enregistrées pour assurer la traçabilité des actions
    \item \textbf{Hashage des mots de passe :} Les mots de passe des utilisateurs sont hashés (BCrypt ou similaire) avant stockage en base de données
    \item \textbf{Sessions :} Gestion des sessions utilisateur avec timeout automatique
\end{itemize}

\subsection{Technologies Utilisées}

\begin{itemize}
    \item \textbf{Backend :} Java EE (Servlets, JSP), PostgreSQL
    \item \textbf{Frontend Web :} JSP, HTML, Tailwind CSS
    \item \textbf{Architecture :} Pattern MVC
    \item \textbf{Serveur d'Application :} Apache Tomcat
    \item \textbf{Conteneurisation :} Docker, Docker Compose
    \item \textbf{Notifications :} Email (SMTP), WhatsApp (API WhatsApp Business ou services tiers)
\end{itemize}

\section{Conclusion}

Ce cahier des charges version 2.0 pose les bases de développement d'un système de gestion des salles pour l'école ENSAA, visant à optimiser l'utilisation des espaces et simplifier le processus de réservation et de gestion des salles.

\textbf{Évolutions majeures de cette version :}
\begin{itemize}
    \item \textbf{Architecture unifiée :} Adoption d'une table unique \texttt{users} avec système de rôles, simplifiant grandement l'architecture et la maintenance
    \item \textbf{Workflow centralisé d'inscription :} Toutes les inscriptions passent par l'approbation du Membre d'Administration, garantissant un meilleur contrôle
    \item \textbf{Rôle Membre d'Administration :} Nouveau rôle central remplaçant le "Responsable des Salles" avec des permissions étendues
    \item \textbf{Gestion des cycles :} Support explicite des cycles Préparatoire (2 ans) et Ingénieur (3 ans) avec gestion par année
    \item \textbf{Groupes TP/TD :} Gestion fine des groupes (Groupe 1 et Groupe 2) pour les travaux pratiques et dirigés
    \item \textbf{Notifications doubles :} Systématisation des notifications Email + WhatsApp pour tous les événements critiques
\end{itemize}

Le système est conçu pour offrir une expérience fluide aux utilisateurs (membre d'administration, coordinateurs, professeurs, et membres de clubs) tout en répondant aux exigences de sécurité, de traçabilité et de simplicité architecturale. L'approche unifiée de la gestion des utilisateurs et le workflow d'approbation centralisé constituent des améliorations majeures par rapport à l'architecture précédente.

\end{document}